\documentclass[11pt]{article}
\usepackage[a4paper,margin=1in]{geometry}
\usepackage{booktabs}
\usepackage{array}
\usepackage{longtable}
\usepackage{enumitem}
\usepackage{hyperref}
\usepackage{xcolor}
\usepackage{listings}
\usepackage{amsmath}
\usepackage{titlesec}

\setlist[itemize]{leftmargin=1.2em,itemsep=2pt,topsep=2pt}
\setlist[enumerate]{leftmargin=1.4em,itemsep=2pt,topsep=2pt}

\lstdefinestyle{py}{
  language=Python,
  basicstyle=\ttfamily\small,
  keywordstyle=\color{blue!65!black},
  commentstyle=\color{green!35!black},
  stringstyle=\color{orange!45!black},
  showstringspaces=false,
  frame=single,
  breaklines=true,
  columns=fullflexible,
  keepspaces=true
}

\titleformat{\section}{\large\bfseries}{}{0em}{}
\titleformat{\subsection}{\normalsize\bfseries}{}{0em}{}

\title{AR525 MC-PILOT Mentor Meeting Crash Brief}
\author{Prepared for quick handover and Q\&A defense}
\date{April 25, 2026}

\begin{document}
\maketitle
\thispagestyle{empty}

\section*{30-second opening script}
We recreated the core MC-PILOT idea on top of MC-PILCO in a pure Python ballistic simulator, then stress-tested it across release-height configurations. The main contribution is not just running baseline code, but diagnosing and fixing four non-obvious training failure modes: broken gradient path, post-landing rollout corruption, cost-scale mismatch, and horizon/indexing bugs. After fixes, baseline and elevated settings converge reliably, and we documented when and why they fail.

\section*{1. What this project is}
\begin{itemize}
\item \textbf{Task:} Throw a ball to targets using data-efficient model-based RL.
\item \textbf{Base algorithm:} MC-PILCO style GP dynamics + Monte Carlo policy optimization.
\item \textbf{Throwing adaptation:} Single-shot policy $\pi(P) \rightarrow$ release speed at $t=0$; free flight afterward.
\item \textbf{State used in baseline:} $\tilde{x} = [x,y,z,v_x,v_y,v_z,P_x,P_y]$ (8-D).
\item \textbf{Main repo reality (current branch):}
  \begin{itemize}
  \item Present: \texttt{mc-pilot/}, \texttt{mc-pilot-elevated/}, \texttt{MC-PILCO/}
  \item Not standalone here (only report-described): \texttt{mc-pilot-pybullet/}, \texttt{mc-pilot-pb-elevated/}
  \end{itemize}
\end{itemize}

\section*{2. Code map you should memorize}
\begin{longtable}{p{0.30\linewidth} p{0.66\linewidth}}
\toprule
\textbf{File} & \textbf{Role} \\
\midrule
\texttt{mc-pilot/test\_mc\_pilot.py} & Baseline training script, hyperparameters, reinforce loop setup \\
\texttt{mc-pilot/policy\_learning/Policy.py} & Throwing policy and exploration policies \\
\texttt{mc-pilot/model\_learning/Model\_learning.py} & GP model class for ballistic dynamics \\
\texttt{mc-pilot/policy\_learning/MC\_PILCO.py} & Core rollout + policy optimization logic (key fixes live here) \\
\texttt{mc-pilot/simulation\_class/model.py} & Ballistic simulator and drag model \\
\texttt{mc-pilot-elevated/test\_mc\_pilot\_\{b,c,d,e\}.py} & Release-height experiments \\
\texttt{mc-pilot-elevated/policy\_learning/Policy.py} & Stratified exploration implementation \\
\bottomrule
\end{longtable}

\section*{3. Pipeline in one minute}
\subsection*{3.1 Training loop pseudocode}
\begin{lstlisting}[style=py]
for trial in range(num_trials):
    # 1) collect throw data (exploration policy)
    for k in range(Nexp if trial == 0 else 1):
        target = sample_target()
        s0 = [release_pos, zero_vel, target]
        trajectory = ThrowingSystem.rollout(s0, exploration_policy)
        GP.add_data(trajectory)

    # 2) fit GP dynamics model on collected transitions
    GP.optimize_marginal_likelihood()

    # 3) optimize policy with Monte Carlo particles
    for step in range(Nopt):
        targets = sample_targets(M)
        speed = policy(targets)  # one-shot speed at t=0
        v0 = speed_to_3d_velocity(speed, targets)
        particles = init_particles(release_pos, v0, targets)

        particles = rollout_with_GP(particles, freeze_if_landed=True)
        cost = landing_cost(particles[-1], targets, lc)
        backprop(cost)
\end{lstlisting}

\subsection*{3.2 Core equations}
\textbf{Policy (RBF + squash):}
\[
\pi_{\theta}(P)=\frac{u_M}{2}\left(\tanh\!\left(\sum_{i=1}^{N_b}w_i\exp\!\left(-\frac{\|a_i-P\|^2}{\ell_s^2}\right)\right)+1\right)
\]

\textbf{Speed to release velocity:}
\[
\phi=\operatorname{atan2}(P_y-y_0,\,P_x-x_0),\quad
v_x=u\cos\alpha\cos\phi,\;v_y=u\cos\alpha\sin\phi,\;v_z=u\sin\alpha
\]

\textbf{Landing cost:}
\[
c(\tilde{x}_T)=1-\exp\!\left(-\frac{\|p_T-P\|^2}{\ell_c^2}\right)
\]

\textbf{State update from GP velocity deltas:}
\[
v_{t+1}=v_t+\Delta v_t,\qquad
p_{t+1}=p_t+T_s v_t+\frac{T_s}{2}\Delta v_t
\]

\textbf{Drag with air-relative velocity (wind-aware physics):}
\[
F_D=-\tfrac{1}{2}\rho C_D A\|v_{rel}\|v_{rel},\quad v_{rel}=v-w
\]

\section*{4. Important code snippets (real project snippets)}
\subsection*{4.1 Baseline hyperparameters and horizon fix}
\begin{lstlisting}[style=py]
# mc-pilot/test_mc_pilot.py
Nexp = 5
Nopt = 1500
M    = 400
Nb   = 250
uM   = 3.5
Ts   = 0.02
T    = 0.7
lc   = 0.5

n_list = Nexp + num_trials
policy_optimization_dict["opt_steps_list"] = [Nopt] * n_list
reinforce_param_dict["T_control"] = T
\end{lstlisting}

\subsection*{4.2 Gradient-preserving one-shot apply\_policy}
\begin{lstlisting}[style=py]
# mc-pilot/policy_learning/MC_PILCO.py (MC_PILOT.apply_policy)
speed = self.control_policy(policy_input, t=0, p_dropout=p_dropout)
azimuth = torch.atan2(dy, dx)
vx = speed * torch.cos(alpha) * torch.cos(azimuth)
vy = speed * torch.cos(alpha) * torch.sin(azimuth)
vz = speed * torch.sin(alpha)
v3d = torch.cat([vx, vy, vz], dim=1)
init_particles = torch.cat([ball_pos, v3d, targets], dim=1)
\end{lstlisting}

\subsection*{4.3 Landing freeze logic}
\begin{lstlisting}[style=py]
landed = torch.zeros(num_particles, dtype=torch.bool, device=self.device)
for t in range(1, int(T_control)):
    next_particles, _, _ = self.model_learning.get_next_state(...)
    prev = states_sequence_list[t - 1]
    frozen = torch.where(landed.unsqueeze(1), prev, next_particles)
    landed = landed | (next_particles[:, 2] <= 0.0)
    states_sequence_list.append(frozen)
\end{lstlisting}

\subsection*{4.4 Stratified exploration}
\begin{lstlisting}[style=py]
# mc-pilot-elevated/policy_learning/Policy.py
stratum = self._throw_idx % self.n_strata
band = (self.u_max - self.u_min) / self.n_strata
lo = self.u_min + stratum * band
hi = self.u_min + (stratum + 1) * band
speed = lo + (hi - lo) * torch.rand(1, dtype=self.dtype, device=self.device)
self._throw_idx += 1
\end{lstlisting}

\subsection*{4.5 Air-relative drag}
\begin{lstlisting}[style=py]
# mc-pilot/simulation_class/model.py
v_rel = vel - wind
v_rel_norm = np.linalg.norm(v_rel)
F_drag = -0.5 * _RHO * cd * A * v_rel_norm * v_rel
\end{lstlisting}

\section*{5. Four failure modes and fixes (high-value Q\&A material)}
\begin{longtable}{p{0.23\linewidth} p{0.30\linewidth} p{0.39\linewidth}}
\toprule
\textbf{Failure mode} & \textbf{What went wrong} & \textbf{Fix and file} \\
\midrule
Gradient path broken & Policy speed was not effectively connected to GP rollout chain & Embed speed-derived $v_{3d}$ directly into initial particle state in \texttt{mc-pilot/policy\_learning/MC\_PILCO.py} \\
Post-landing hallucination & GP kept propagating particles underground, corrupting terminal cost & Add landed mask + freeze logic in \texttt{mc-pilot/policy\_learning/MC\_PILCO.py} \\
Cost saturation & $\ell_c=0.1$ gave near-zero gradients when early errors were large & Use $\ell_c=0.5$ for bootstrap phase in \texttt{mc-pilot/test\_mc\_pilot.py} (and adjusted variants in elevated scripts) \\
Horizon/indexing bugs & Wrong interpretation of control horizon and too-short opt list indexing & Set \texttt{T\_control=T} and size lists as \texttt{Nexp + num\_trials} in \texttt{mc-pilot/test\_mc\_pilot.py} \\
\bottomrule
\end{longtable}

\section*{6. Results you can quote quickly}
\subsection*{6.1 Baseline (mc-pilot)}
\begin{itemize}
\item Current best baseline: 5/5 policy throws hit in seed=1 run after exploration phase.
\item Typical trial errors: 85 mm, 33 mm, 23 mm, 17 mm, 8 mm (improves over trials).
\item Key working setup: $T_s=0.02$, $T=0.7$, $\ell_c=0.5$, $l_M=1.75$.
\end{itemize}

\subsection*{6.2 Elevated release study (mc-pilot-elevated)}
\begin{longtable}{p{0.18\linewidth} p{0.13\linewidth} p{0.16\linewidth} p{0.16\linewidth} p{0.29\linewidth}}
\toprule
\textbf{Config} & \textbf{$z_{release}$} & \textbf{Exploration} & \textbf{Result} & \textbf{Takeaway} \\
\midrule
B random (seed=1) & 1.0 m & random & 0/10 & bad speed coverage can fail completely \\
B-strat & 1.0 m & stratified & 10/10 & seed-robust after coverage fix \\
C-strat & 1.5 m & stratified & 10/10 & mean error about 14 mm \\
D-strat & 2.0 m & stratified & 10/10 & mean error about 16 mm \\
E-strat1 to E-strat4 & 0.0 m & multiple attempts & 0/10 & narrow target range broke policy resolution \\
E-strat5 & 0.0 m & targeted + tuned & 10/10 & fixed by shrinking policy lengthscale to 0.15 \\
\bottomrule
\end{longtable}

\section*{7. Mentor-safe boundary statements}
\begin{itemize}
\item In this branch, pybullet helpers exist, but standalone folders \texttt{mc-pilot-pybullet/} and \texttt{mc-pilot-pb-elevated/} are report-described, not present as separate trees.
\item Wind physics (air-relative drag) exists in simulator code; full 8-D to 11-D policy and GP wiring described in report is not fully visible in this checkout.
\item Current strongest claims are simulation-based and geometry-generalization-focused, not hardware deployment claims.
\end{itemize}

\section*{8. Rapid-fire questions and compact answers}
\begin{enumerate}
\item \textbf{What is novel here beyond running MC-PILCO?} \\
We adapted it to single-shot throwing and diagnosed four hidden failure modes with validated fixes.

\item \textbf{Why not use paper defaults directly?} \\
Because geometry and simulator details changed gradient behavior; some defaults caused saturation or instability.

\item \textbf{Why is stratified exploration important?} \\
With only 5 exploration throws, random seeds can miss crucial speed bands; stratified guarantees coverage.

\item \textbf{Why did Config E fail repeatedly?} \\
Target range was narrow, but policy lengthscale stayed large, making RBF activations nearly identical across targets.

\item \textbf{What fixed Config E finally?} \\
Reducing policy lengthscale to 0.15, so targets became distinguishable and policy stopped outputting near-constant max speed.

\item \textbf{Why use $\ell_c=0.5$ instead of 0.1?} \\
0.1 saturated cost early and killed gradients; 0.5 preserved gradient signal during bootstrap.

\item \textbf{Where is the core learning logic?} \\
In \texttt{mc-pilot/policy\_learning/MC\_PILCO.py}, especially \texttt{MC\_PILOT.apply\_policy()}.

\item \textbf{How is wind handled currently?} \\
Through air-relative drag in simulator physics: $v_{rel}=v-w$.

\item \textbf{What is the strongest current result statement?} \\
After fixes, baseline converges reliably and elevated B/C/D plus E-strat5 achieve perfect trial hit-rates in tested seeds.

\item \textbf{What are the next serious steps?} \\
Multi-seed statistics, cleaner wind-aware full-state wiring, and physical robot validation.
\end{enumerate}

\section*{9. Reproduce quickly (commands)}
\begin{lstlisting}[style=py]
# Baseline
python mc-pilot/test_mc_pilot.py -seed 1 -num_trials 10

# Elevated examples
python mc-pilot-elevated/test_mc_pilot_b_strat.py -seed 1 -num_trials 10
python mc-pilot-elevated/test_mc_pilot_e_strat5.py -seed 1 -num_trials 10
\end{lstlisting}

\section*{10. One-line closing you can use}
This project is strong because it combines algorithm adaptation, failure diagnosis, and reproducible improvements, with honest boundaries on what is fully implemented in this branch.

\end{document}
